
%\addcontentsline{toc}{section}{В. \,Реалізація проєкту Периметр-Фортеця }

% \addtocontents{toc}{\protect\setcounter{tocdepth}{-1}}

\section {\centering Реалізація проєкту ``Периметр-Фортеця''}

Проєктна робота ``Периметр-Фортеця''  присвячена архітектурному та політичному проєктуванню комплексної системи інформаційної безпеки для малої корпоративної мережі (МКМ) на прикладі юридичної компанії. Основною метою було досягнення оптимального балансу тріади Конфіденційність, Цілісність, Доступність при мінімізації економічних витрат. Методологія реалізації базувалася на застосуванні парадигми ``Заборонити за замовчуванням'' (Deny by Default) на рівні мережевого периметру та впровадженні Дискреційної Моделі Керування Доступом (DAC) на рівні внутрішнього файлового сервера. У результаті розроблено логічну конфігурацію міжмережевого екрану, що забезпечує захист від зовнішніх DoS-атак та несанкціонованого сканування портів, а також політику безпечного віддаленого доступу через VPN-тунель (OpenVPN). Проєкт демонструє інтеграцію мережевих, системних та організаційних засобів інформаційної безпеки для нейтралізації інсайдерських та зовнішніх загроз.

Мета такого проєкту полягає у обґрунтуванні, розробці та документуванні архітектури кіберзахисту малої корпоративної мережі, що забезпечує стійкість до типових зовнішніх та внутрішніх загроз, з особливим акцентом на захисті конфіденційної інформації та безпечному віддаленому доступі.

До проєкту можна поставити таке проблемне питання - Як спроєктувати економічно ефективний, але надійний мережевий периметр для малої компанії, який забезпечить конфіденційність конфіденційних даних (тріада CIA), одночасно надаючи безпечний віддалений доступ, і чи може брандмауер повністю захистити від загроз, що виникають зсередини мережі (інсайдерські ризики)?

Ключові питання до цього проєкту можна запропонувати наступні:
\begin{itemize}
    \item Які базові правила фільтрації (MAC, IP, порт) є критично необхідними для захисту внутрішнього сервера від DoS-атак?
    \item Як обрана модель керування доступом на сервері узгоджується з загальною політикою безпеки мережі?
    \item Як брандмауер та політика безпеки забезпечують баланс між конфіденційністю, цілісністю та доступністю (тріада CIA) даних у цій мережі?
\end{itemize}

Реалізація проєкту ґрунтується на двовекторному підході до захисту:
\begin{enumerate}
    \item Периметровий захист - встановлено міжмережевий екран, який працює у режимі фільтрації пакетів. Ядром конфігурації є правило (DROP ALL), що гарантує блокування будь-якого вхідного (Inbound) трафіку, крім:
    \begin{itemize}
        \item Захищеного віддаленого доступу (VPN): Дозволено лише трафік через порт 1194 (OpenVPN) від авторизованого WAN-адреса.
        \item Захисту Доступності: Встановлено обмеження на ICMP-трафік (Правило 1) для запобігання DoS-атакам типу ``ping-флу''.
    \end{itemize}
    \item Внутрішній системний захист - застосування дискреційної моделі керування доступом (DAC) на файловому сервері (192.168.1.10). Ця модель мінімізує інсайдерський ризик, оскільки власники даних самостійно контролюють права ``читання/запису/виконання'', що є ефективним доповненням до мережевого захисту, який не здатен контролювати дії авторизованих внутрішніх користувачів.
    \item Організаційний захист - введено політику стратегії 3-2-1 для резервного копіювання, що забезпечує високу Цілісність та Доступність критичних даних навіть у разі апаратних збоїв чи успішного кіберінциденту.
\end{enumerate}

У випадку виконання проєкту групою учнів, рекомендується розподіл функціональних ролей (таблиця~\ref{RolesPproject}), що забезпечує глибоке опрацювання всіх розділів модуля.

\begin{longtable}{|
    >{\raggedright\arraybackslash}p{0.18\linewidth}| % Ролі
    >{\raggedright\arraybackslash}p{0.3\linewidth}| % Основна сфера
    >{\raggedright\arraybackslash}p{0.45\linewidth}| % Ключові завдання
    }
    
    % --- ЗАГОЛОВОК ТАБЛИЦІ ---
    \caption{Розподіл ролей у проєкті}
    \label{RolesPproject} \\

    % --- ШАПКА ПЕРШОЇ СТОРІНКИ ---
    \hline
    \textbf{Ролі} & \textbf{Основна сфера відповідальності} & \textbf{Ключові завдання у проєкті} \\
    \hline
    \endfirsthead

    \multicolumn{3}{r}{{\normalsize Продовження табл. \thetable}} \\
    \hline

    \hline
    \textbf{Ролі} & \textbf{Основна сфера відповідальності} & \textbf{Ключові завдання у проєкті} \\
    \hline
    \endhead

    % --- НИЖНЯ ЛІНІЯ ---
    \hline
    \endlastfoot

    % --- РЯДОК 1 ---
    Аналітик ризиків та політик & 
    Організаційні та правові основи ІБ. & 
    \begin{minipage}[t]{\linewidth}
    \vspace{1mm}
    Проведення аналізу загроз, формулювання \textbf{``Документу політики безпеки''}.
    \vspace{1mm}
    \end{minipage} \\
    \hline

    % --- РЯДОК 2 ---
    Архітектор мережі та периметру & 
    Мережевий захист та логічна структура. & 
    \begin{minipage}[t]{\linewidth}
    \vspace{1mm}
    Розробка схеми мережі, обґрунтування принципів роботи брандмауера, розробка файлу конфігурації брандмауера.
    \vspace{1mm}
    \end{minipage} \\
    \hline

    % --- РЯДОК 3 ---
    Інженер систем безпеки & 
    Системний захист та керування доступом. & 
    \begin{minipage}[t]{\linewidth}
    \vspace{1mm}
    Обґрунтування моделі DAC, визначення прав доступу на сервері, розробка Політики Резервного Копіювання.
    \vspace{1mm}
    \end{minipage} \\
    \hline

    % --- РЯДОК 4 ---
    Технічний документатор та презентатор & 
    Фінальна звітність та підготовка до захисту. & 
    Перевірка узгодженості розділів, підготовка презентації. \\

\end{longtable}


При розробці проєкту учні зосереджуються на різних питаннях розділів вибіркового модуля. У контексті інформаційної безпеки, першочергова увага має приділятися об'єктам захисту (Розділ І). Найбільш критичними активами є внутрішній сервер, що зберігає конфіденційні дані клієнтів, та робочі станції користувачів, які залишаються вразливими через ризики соціальної інженерії та загрози шкідливого програмного забезпечення. Типові загрози включають фішинг, використання відомих вразливостей застарілого програмного забезпечення, а також зовнішні мережеві атаки, як-от сканування портів чи DDoS/DoS-атаки, спрямовані на периметр. Зосередженість на таких питаннях дозволяє учням чітко усвідомити, що захист периметра (брандмауер) та внутрішній контроль (керування доступом) є однаково критичними елементами цілісної стратегії інформаційної безпеки.

Переходячи до системного захисту (Розділ ІІ), ключовим кроком є обґрунтування відповідної моделі керування доступом. Для внутрішнього сервера, що є сховищем чутливих даних, оптимально обрано дискреційну модель керування доступом (DAC). Її реалізація передбачає, що користувачі мають чітко визначені права доступу (читання/запис) лише до власних тек або тек, до яких вони явно отримали дозвіл. Керівник або адміністратор системи, відповідно, зберігає повний контроль над ресурсами. Така модель узгоджується з політикою гнучкості, необхідною для динамічної роботи, водночас покладаючи відповідальність за призначення прав доступу безпосередньо на власників ресурсів.

На рівні мережевого захисту (Розділ ІІІ), основним елементом є брандмауер, який функціонує за фундаментальним принципом ``Заборонити все, що явно не дозволено'' (Deny by Default). Це забезпечує максимальну безпеку, блокуючи весь потенційно небезпечний трафік за замовчуванням. Крім того, застосування NAT (трансляція мережевих адрес) є необхідним для приховування внутрішніх IP-адрес мережі від зовнішнього світу, що значно ускладнює прямі атаки. Таким чином, брандмауер виступає як єдина контрольована точка входу та виходу, дозволяючи моніторити трафік та ефективно запобігати несанкціонованому скануванню та доступу ззовні.

Результатом проєкту  є схема мережі, документ політики безпеки ``ЮрЗахист'', файл конфігурації брандмауера .

Схема мережі, яка спроєктована за топологією ``Зірка'' з єдиною точкою виходу — міжмережевим екраном. Компонентами такої мережі є 
\begin{itemize}
    \item міжмережевий екран (firewall/router): зовнішній IP (WAN) та внутрішній IP (LAN: 192.168.1.1)
    \item внутрішній сервер: статичний IP 192.168.1.10 (зберігання даних)
    \item робочі станції (x3): динамічні IP (192.168.1.50-52)
\end{itemize}
Документ політики безпеки ``ЮрЗахист'' складається з розділів що подані у таблиці~\ref{tab:appendv_yur}

% Requires: \usepackage{array}
\begin{table}[h!]
    \caption{Розділи політики безпеки}
    \label{tab:appendv_yur}
    \begin{tabular}{|>{\raggedright\arraybackslash}p{5.5cm}|>{\raggedright\arraybackslash}p{9.5cm}|}
        \hline
        \textbf{Назва політики} & \textbf{Визначення та Реалізація} \\
        \hline
        Політика керування доступом & На сервері використовується модель DAC. Усі користувачі повинні мати унікальні, складні паролі (мін. 12 символів, 4 типи). Після 5 невдалих спроб входу обліковий запис блокується на 1 годину. \\
        \hline
        Політика резервного копіювання (Цілісність/Доступність) & Повне щоденне резервне копіювання даних із сервера о 23:00. Копії зберігаються за моделлю 3-2-1 (3 копії, на 2 різних носіях, 1 копія — поза офісом). \\
        \hline
        Політика мережевого доступу & Заборонений будь-який вхідний (Inbound) трафік, окрім спеціально дозволеного VPN-з'єднання для керівника. Увесь інший доступ до внутрішньої мережі блокується на рівні брандмауера. \\
        \hline
        Політика оновлень & Усе програмне забезпечення на сервері та робочих станціях має бути оновлено до останньої версії протягом 7 днів після виходу патча. \\
        \hline
    \end{tabular}
\end{table}

Файл Конфігурації Брандмауера будується за логікою правил. Налаштування виконується на міжмережевому екрані, що захищає периметр мережі. Правила перевіряються по черзі (зверху донизу), і якщо правило спрацьовує, подальша перевірка не відбувається.

% Requires: \usepackage{array}
\begin{table}[h]
    \small
    \centering
    \caption{Правила для Firewall}
    \label{tab:appendix_firewall}
    \begin{tabular}{|p{1.8cm}|p{1.8cm}|p{1.5cm}|p{2.4cm}|p{2cm}|p{2cm}|p{2.8cm}|}
        \hline
        Пріоритет & Дія & Прото-кол & Джерело (Source) & Призна-чення (Destination) & Порт & Мета та обґрунтування \\
        \hline
        1 (Анти-DoS) & DROP & ICMP & ANY & 192.168.1.1 (Firewall) & ANY & Обмеження ICMP-пакетів для захисту від пінг-флуду (DoS-атаки). \\
        \hline
        2 (Відда-лений доступ) & ACCEPT & TCP/UDP & WAN (IP керівника) & 192.168.1.1 (Firewall) & 1194 (OpenVPN) & Дозвіл VPN-з'єднання для безпечного віддаленого доступу керівника. \\
        \hline
        3 (Вихідний трафік) & ACCEPT & TCP & 192.168.1.0/24 (LAN) & ANY & 80, 443, 53 & Дозвіл внутрішнім користувачам доступу до Інтернету (HTTP, HTTPS, DNS). \\
        \hline
        4 (Внутріш-ній доступ) & ACCEPT & ANY & 192.168.1.0/24 (LAN) & 192.168.1.10 (Server) & ANY & Дозвіл доступу робочих станцій до сервера (контролюється DAC на рівні ОС сервера). \\
        \hline
        5 (Заборонити) & DROP & ANY & ANY & ANY & ANY & Правило ``Заборонити все, що не дозволено'' (Default Deny). Блокує всі інші пакети. \\
        \hline
    \end{tabular}
\end{table}



Таким чином, інтеграція правил міжмережевого екрану, політики безпеки та керування доступом на системному рівні створює надійну ``Фортецю'' для даних компанії.